\section{课程安排与核心内容}

\subsection{学习目标与技术范围}
CUDA 存储模型连接了两类问题：线程能够访问哪些数据，以及这些访问如何影响计算速度。基本矩阵乘法提供了一个完整例子：为输出元素分配线程，建立正确的地址映射，再根据浮点运算量与访存量分析性能。

学习目标为：了解 CUDA 线程可访问的各类存储器的基本特点；为 Lab 2 的基本矩阵乘法实现做准备；学习评估全局内存访问对性能的影响；了解 GPU 性能指标。对应的技术内容依次为存储层次、基本矩阵乘法，以及算术强度和 Roofline 性能模型。

\subsection{实验、交流与课程项目}
\subsubsection{实验与支持渠道}
以下安排对应 \textbf{2026 年 9 月 8 日}的课程信息。Lab 1 的截止时间为当周周五，\textbf{9 月 11 日}；需要先完成 Lab 0，Lab 1 不计入最终成绩。

课程问题通过 Campuswire 发布，课程工作人员只回复在 Campuswire 上发布的问题。提问前应先尝试自行查找答案，包括询问课程聊天机器人。答疑时间同时包含线下与线上形式，需要查看相应时间表。课后事项为阅读教材第 3 章、取得 Delta 账户、设置 GitHub，以及提交 Labs 0 和 1。

\subsubsection{课程项目与暂定里程碑}
两个项目方向都只实现\term{推理的前向计算}{inference / forward pass}，区别在于实现对象与组织形式。
\begin{table}[H]
\centering\small
\begin{tabularx}{\textwidth}{L{0.12\textwidth} Y L{0.25\textwidth}}
\toprule
方向 & 实现内容 & 组织形式 \\
\midrule
CNN & LeNet-5 卷积神经网络中的卷积层 & 个人项目 \\
GPT-2 & GPT-2 模型中的各层 & 团队项目，每组 3--4 人 \\
\bottomrule
\end{tabularx}
\end{table}
\noindent 暂定时间为：Milestone 1 在第 7--8 周，Milestone 2 在第 10--11 周，Milestone 3 在第 15--16 周。以上里程碑为暂定周次。
